\documentclass[11pt]{article}
\usepackage[margin=1in]{geometry}
\usepackage{amsmath,amssymb,amsthm}
\usepackage[T1]{fontenc}
\usepackage{lmodern}
\usepackage{microtype}
\usepackage[hidelinks]{hyperref}
\usepackage{xurl}
\newtheorem{theorem}{Theorem}
\newtheorem{lemma}{Lemma}
\newcommand{\R}{\mathbb R}
\newcommand{\C}{\mathbb C}
\newcommand{\Q}{\mathbb Q}
\newcommand{\Z}{\mathbb Z}
\newcommand{\code}[1]{\nolinkurl{#1}}
\title{Conjecture 00000000374:\\zero is excluded at the claimed full-line endpoint}
\author{C0ldSmi1e}
\date{October 9, 2026}
\begin{document}
\maketitle
\begin{abstract}
The conjecture asserts that every real number admits infinitely many distinct
algebraic approximants of degree at most two with error less than
$H(\alpha)^{-\tau}$ when $\tau\leq3$. At $\tau=3$, zero does not have this property.
We prove this for each of three explicit standard arithmetic height conventions:
naive coefficient height, Mahler measure, and absolute multiplicative Weil
height. The proof allows either real or complex algebraic approximants.
A Lean 4 formalization proves the negation of the full-real-line conjunct with
the stated degree, strict inequality, exponent, and infinitude quantifiers.
\end{abstract}

\section{Statement and scope}
In both supplied language versions, $P(\tau)$ consists of real $\xi$ for which
infinitely many algebraic numbers $\alpha$ of degree at most two satisfy
$|\xi-\alpha|<H(\alpha)^{-\tau}$. The conjecture combines three assertions:
Hausdorff dimension $3/\tau$ for $\tau>3$; $P(\tau)=\R$ for $\tau\leq3$;
and continuity of the dimension formula at $3$. We disprove the second assertion,
so their conjunction is false. We make no separate claim about the other two.
The original does not exclude rational targets, replace equality with an
almost-everywhere assertion, or permit repeated copies of one approximant.
The counterexample uses the included endpoint $\tau=3$ and the real target $0$.
The exact bilingual source accompanies this submission as \code{ORIGINAL.md}.

The source does not define $H$. Accordingly, this report specifies and treats
three conventional multiplicative arithmetic heights separately, without
identifying their values or changing the exponent. Our result covers those
three readings; it is not a theorem about an arbitrary function denoted by $H$.
Both real and complex approximant domains are covered because the source leaves
that choice implicit. Absolute value on $\C$ means its usual modulus.

\section{Actual algebraic numbers and their heights}
Let $\alpha\in\C$ be algebraic over $\Q$, with degree $d$. Let
\[
 p_\alpha(X)=a_dX^d+\cdots+a_0\in\Z[X]
\]
be its primitive irreducible integer minimal polynomial, normalized by $a_d>0$.
If its complex roots, with multiplicities, are $\alpha_1,\ldots,\alpha_d$, define
\begin{align*}
 H_N(\alpha)&=\max_{0\leq i\leq d}|a_i|,\\
 M(\alpha)&=|a_d|\prod_{i=1}^d\max(1,|\alpha_i|),\\
 W(\alpha)&=M(\alpha)^{1/d}.
\end{align*}
The last expression is the primitive-minimal-polynomial definition of absolute
multiplicative Weil height. Its degree normalization uses the actual $d$, not
the upper bound $2$: $W^d=M$. This formalization does not separately construct
the equivalent product over places of a number field.

The Lean construction starts with the monic rational minimal polynomial,
clears denominators, and takes the primitive part. It proves a nonzero rational
scalar relation to the monic polynomial, the same degree, irreducibility,
vanishing at $\alpha$, and a nonzero constant coefficient if $\alpha\ne0$.
Gauss's lemma gives uniqueness up to sign. A positive-leading-coefficient
representative is constructed and proved unique. Each height is sign invariant;
therefore its implemented value agrees exactly with the displayed formula for
any primitive irreducible integer polynomial vanishing at $\alpha$.
This includes algebraic numbers that are not algebraic integers.

\section{Finitely many approximants to zero}
For a chosen height $H$, put
\[
 A_H=\{\alpha\in\C:\alpha\text{ is algebraic over }\Q,\quad
 [\Q(\alpha):\Q]\leq2,\quad |\alpha|<H(\alpha)^{-3}\}.
\]
We prove $A_H$ finite. Restricting to real approximants can only decrease it.

\begin{lemma}[Naive height]
$A_{H_N}$ is finite.
\end{lemma}
\begin{proof}
For $\alpha\ne0$, write $r=|\alpha|$ and $H=H_N(\alpha)$.
The integer $a_0$ is nonzero. Since $d\leq2$, padding missing coefficients
with zeros gives
\[
 1\leq |a_0|=|a_1\alpha+a_2\alpha^2|\leq Hr+Hr^2.
\]
If $H\geq2$ and $r<H^{-3}$, then $rH^3<1$. Since $H^2\geq2$,
we obtain $2Hr<1$. Also $0\leq r\leq1$, so
$Hr+Hr^2\leq2Hr<1$, a contradiction.
Every qualifying nonzero $\alpha$ therefore has $H\leq1$.
Its minimal polynomial has degree at most two and all coefficients in
$\{-1,0,1\}$. There are finitely many such polynomials and each nonzero
polynomial has finitely many roots. Their finite union, together with $0$,
contains $A_{H_N}$. This is a symbolic finiteness argument, not a bounded
numerical search.
\end{proof}

\begin{lemma}[Mahler measure and Weil height]
$A_M\subseteq\{0\}$ and $A_W\subseteq\{0\}$.
\end{lemma}
\begin{proof}
Let $\alpha\ne0$ have degree $d\leq2$. The constant-term product formula yields
\[
 1\leq |a_0|=|a_d|\prod_{i=1}^d|\alpha_i|\leq |\alpha|M(\alpha).
\]
For the last inequality, separate one occurrence of the root $\alpha$, bound
every remaining modulus by its maximum with $1$, and, if necessary, multiply
by the extra factor $\max(1,|\alpha|)\geq1$.
Since $|a_d|\geq1$, we have $M\geq1$ and hence $W\geq1$.
The actual algebraic degree is $1$ or $2$, so $M=W^d\leq W^2$.
Consequently
\[
 1\leq |\alpha|M^2,\qquad 1\leq |\alpha|W^2.
\]
For either $H=M$ or $H=W$, a strict inequality
$|\alpha|<H^{-3}\leq H^{-2}$ would imply $|\alpha|H^2<1$,
contradicting these bounds. Thus no nonzero $\alpha$ qualifies.
\end{proof}

\begin{theorem}
For each $H\in\{H_N,M,W\}$ and either real or complex algebraic approximants,
$0\notin P(3)$. In particular,
\[
 \neg\bigl(\forall\tau\in\R,\ \tau\leq3\Longrightarrow P(\tau)=\R\bigr).
\]
\end{theorem}
\begin{proof}
The two lemmas show that the set of eligible approximants to $0$ is finite,
and hence cannot be infinite. For real approximants, take its preimage under
the injective map $\R\hookrightarrow\C$. The norm of the difference becomes
the real absolute value, and the rational minimal polynomial is unchanged
under this embedding. If the displayed universal assertion held, its
$\tau=3$ instance would put $0$ in $P(3)$, a contradiction.
\end{proof}

\section{Formal statement and verification boundary}
\code{DegreeAtMostTwo} requires algebraicity explicitly and bounds the natural
degree of \code{minpoly} over $\Q$. Transcendental numbers cannot enter through
the convention that their minimal polynomial is zero.
\code{approximants} is an actual set of complex numbers, and \code{Pcomplex}
uses \code{Set.Infinite}; \code{Preal} uses an actual set of real numbers.
All parameters $\xi,\tau$ are real. The inequalities use real exponentiation,
the literal constant $1$, and the strict comparison in the source.

The final theorem \code{TLMC374.source_conjunct_false} quantifies over
\code{ApproximationDomain} (real or complex) and \code{HeightConvention}
(naive, Mahler, or absolute Weil). Its conclusion is exactly the negation of
$\forall\tau\leq3,\ P(\tau)=\R$ for the selected definitions.
It follows from \code{TLMC374.zero_not_mem_P_three}.
The formalized set and universal assertion exist independently of the chosen
counterexample; no witness-specific surrogate replaces them.

\code{Counterexample.lean} contains the minimal-polynomial bridges, naive
height argument, approximation sets, real-domain transfer, and final theorem.
\code{HeightRoots.lean} proves the root-product bounds and Mahler/Weil
identities. The project pins Lean 4.19.0 and Mathlib v4.19.0, revision
\code{c44e0c8ee63ca166450922a373c7409c5d26b00b}.
The accompanying verification records distinguish independent local checks
from maintainer acceptance. The supplied runner rebuilds authored modules,
replays source with warnings treated as errors, and audits compiled
declarations in both mathematical modules, including generated declarations.
Safe mathematical declarations may depend only on the standard foundational
axioms \code{propext}, \code{Classical.choice}, and \code{Quot.sound}.
The inventory also records compiler-generated unsafe execution helpers;
none is reachable through the type or value dependencies of any safe
mathematical declaration. No numerical computation is needed for the proof.

The missing definition of $H$ remains a limitation of the original statement.
Logarithmic height, an archimedean house omitting denominator information,
and arbitrary rescalings are different functions and are not covered by this
theorem. No assertion is made that the three heights above are equal or can
be exchanged without affecting other parts of the conjecture.
The report resolves the literal full-line assertion under each of the three
explicit standard arithmetic conventions; it does not resolve a newly
specified arbitrary-height problem or determine the Hausdorff dimension.
\end{document}
